%%% Local Variables:
%%% mode: LaTeX
%%% TeX-engine: luatex
%%% TeX-master: t
%%% TeX-backend: "biber"
%%% End:
\documentclass[article,spanish]{apuntes-scr}

\subject{Manual de Referencia Oficial}
\title{Guía Completa de Uso y Configuración de Kitty}
\subtitle{Emulador de Terminal Acelerado por GPU con Modus Vivendi Deuteranopia}
\author{Campos Mendoza, D. Daniel}
\date{%
    \textit{Ecosistema GNU Emacs 30 / Sway / Wayland} \\[0.3em]
    \today\ \textperiodcentered\ \DTMcurrenttime
}

% =========================================================
% BIBLIOGRAFÍA
% =========================================================
% \addbibresource{references.bib}

\begin{document}

\maketitle
\tableofcontents
\vspace{1.5em}\hrule\vspace{1.5em}

\section{Arquitectura y Características de Kitty}
\textbf{Kitty} es un emulador de terminal moderno diseñado para un rendimiento extremo. Traslada el renderizado de texto y glifos a la tarjeta gráfica mediante OpenGL, soporta el protocolo extendido de teclado de Emacs, renderizado tipográfico complejo y gestión nativa de paneles (\textit{splits}) y pestañas sin sobrecarga de multiplexores como tmux.

\section{Tipografía y Renderizado Matemático}
La terminal utiliza una combinación tipográfica de máxima legibilidad:
\begin{itemize}[leftmargin=*]
    \item \textbf{Fuente Principal}: \texttt{Iosevka Term} (Light, 12.5pt) con altura de línea optimizada al 110\% (\texttt{modify\_font cell\_height 110\%}).
    \item \textbf{Símbolos Matemáticos}: Integración mediante \texttt{symbol\_map} directo a \texttt{STIX Two Math} para caracteres de operadores lógicos ($\forall, \exists, \in, \notin$), flechas ($\implies, \to, \mapsto$) y conjuntos numéricos ($\mathbb{R}, \mathbb{C}, \mathbb{N}, \mathbb{Z}$).
\end{itemize}

\section{Gestión de Paneles Divididos (Splits)}
Tu configuración tiene habilitado el motor \texttt{enabled\_layouts splits,stack,tall,fat} que permite dividir el espacio de trabajo de forma intuitiva:

\subsection{Creación y Cierre de Paneles}
\begin{itemize}[leftmargin=*]
    \item \textbf{\texttt{Ctrl + Shift + E}}: Dividir la ventana actual \textbf{verticalmente} (\textit{vsplit}), abriendo la nueva sesión en el mismo directorio de trabajo (\texttt{cwd=current}).
    \item \textbf{\texttt{Ctrl + Shift + D}}: Dividir la ventana actual \textbf{horizontalmente} (\textit{hsplit}).
    \item \textbf{\texttt{Ctrl + Shift + W}}: Cerrar el panel enfocado.
\end{itemize}

\subsection{Navegación Direccional entre Paneles}
\begin{itemize}[leftmargin=*]
    \item \textbf{\texttt{Ctrl + Shift + h}} o \textbf{\texttt{Ctrl + Shift + $\leftarrow$}}: Mover el foco al panel izquierdo.
    \item \textbf{\texttt{Ctrl + Shift + j}} o \textbf{\texttt{Ctrl + Shift + $\downarrow$}}: Mover el foco al panel inferior.
    \item \textbf{\texttt{Ctrl + Shift + k}} o \textbf{\texttt{Ctrl + Shift + $\uparrow$}}: Mover el foco al panel superior.
    \item \textbf{\texttt{Ctrl + Shift + l}} o \textbf{\texttt{Ctrl + Shift + $\rightarrow$}}: Mover el foco al panel derecho.
\end{itemize}

\section{Pestañas y Control del Portapapeles}
\begin{itemize}[leftmargin=*]
    \item \textbf{\texttt{Ctrl + Shift + T}} / \textbf{\texttt{Ctrl + Shift + Q}}: Crear una nueva pestaña / Cerrar la pestaña activa.
    \item \textbf{\texttt{Ctrl + Shift + C}} / \textbf{\texttt{Ctrl + Shift + V}}: Copiar texto seleccionado / Pegar desde el portapapeles del sistema Wayland.
    \item \textbf{Soporte OSC 52}: Permite que programas remotos o locales (como Emacs en TTY con \texttt{clipetty}) copien al portapapeles del sistema a través de secuencias de escape estándar.
\end{itemize}

\section{Historial y Scrollback de 30,000 Líneas}
Kitty almacena hasta \textbf{30,000 líneas} de historial en memoria RAM:
\begin{itemize}[leftmargin=*]
    \item \textbf{\texttt{Shift + RePág}} / \textbf{\texttt{Shift + AvPág}}: Desplazamiento rápido por el historial.
    \item \textbf{\texttt{Shift + $\uparrow$}} / \textbf{\texttt{Shift + $\downarrow$}}: Desplazamiento línea a línea.
    \item \textbf{\texttt{Ctrl + Shift + H}}: Abre el historial de la terminal en un visor de texto editable para buscar y extraer bloques de logs.
\end{itemize}

\section{Paleta Modus Vivendi Deuteranopia}
Fondo negro absoluto (\texttt{\#000000}) con acentos en cian turquesa (\texttt{\#00d3d0}), oro (\texttt{\#fec43f}), azul zafiro (\texttt{\#2fafff}), menta (\texttt{\#70d770}) y óxido (\texttt{\#ff6f6f}) para máximo confort visual y accesibilidad cromática.

\end{document}
